\section{Towards Generative World Modeling}
\label{sec:world}

\subsection{Dynamic Scenes and 4-D World Representation}

The step from static \nvs{} to world modeling hinges on modelling
\emph{temporal dynamics}: how does a scene evolve over time, and can
the model generate novel views at arbitrary future timestamps?
The dynamic NeRF and 4DGS methods in §\ref{sec:nerf}--\ref{sec:3dgs}
address this for \emph{observed} scenes given dense multi-view video.
A harder and more practically relevant regime is \emph{generalizable}
4-D generation from sparse or single-image input---precisely the regime
inhabited by world models.

Video diffusion models have emerged as the most promising backbone.
ViewCrafter~\citep{yu2024viewcrafter} augments a video diffusion model with
point-cloud renderings along the target camera trajectory, providing a
geometric scaffold that constrains the generation while the diffusion prior
fills in appearance and fine detail.
CAT4D~\citep{wu2024cat4d} generalises the Cat3D framework to the temporal
dimension, jointly modelling viewpoint and time with a multi-view video
diffusion model, then lifting the output to an explicit 4-D representation.
4DiM~\citep{watson2024_4dim} takes a more direct approach, training a
4-D diffusion model that simultaneously denoises a set of frames at
varying timestamps and viewpoints, learning a joint spatiotemporal prior.
SV3D~\citep{voleti2024sv3d} and SV4D~\citep{xie2024sv4d} extend Stable
Video Diffusion to 3-D and 4-D object generation respectively, conditioning
on orbital camera paths to ensure multi-view consistency.
Consistent4D~\citep{jiang2024consistent4d} enforces temporal consistency
across a 4-D cascade, resolving inter-frame Gaussian drift.
DreamGaussian4D~\citep{ren2023dreamgaussian4d} shows that SDS can be applied
to 4DGS, enabling text-driven 4-D generation without any video training data.

A complementary line of work addresses the \emph{monocular} case.
Diffusion Priors for Dynamic NVS~\citep{li2024diffpriordynamic} treats the
diffusion model as a frame-level regulariser during per-scene 4-D optimisation
from a single monocular video, substantially improving novel-view quality in
under-constrained regions.

\subsection{Neural Scene Simulators for Autonomous Driving}

Autonomous driving provides a uniquely demanding test-bed for NVS-as-world-model:
scenes are large-scale, unbounded, and contain independently moving actors
(vehicles, pedestrians) that must be rendered consistently from closed-loop
ego-motion trajectories.
Three architectural strategies have emerged.

\paragraph{NeRF-based compositional simulators.}
MARS~\citep{wu2023mars} separates the scene into a static NeRF background
and per-agent NeRF foregrounds, allowing independent actor manipulation.
UniSim~\citep{yang2023unisim} extends this to a closed-loop sensor simulator
that supports controllable scene editing---removing actors, changing weather,
inserting novel agents---with photorealistic rendering at the ego sensor.
EmerNeRF~\citep{yang2023emernerf} learns the dynamic-static decomposition
in a self-supervised manner by decomposing scene flow from multi-camera video,
without per-actor instance labels.

\paragraph{3DGS-based driving reconstruction.}
Street Gaussians~\citep{yan2024streetgaussian} extends 3DGS to urban scenes
by representing tracked foreground objects with separate Gaussian models and
compositing them over a background Gaussian field.
PVG~\citep{chen2023pvg} models Gaussian motion with periodic vibration functions
that efficiently encode slow, cyclic motions common in urban traffic.
OmniRe~\citep{chen2024omnire} unifies background, foreground, and sky in a
single Gaussian framework, establishing a strong all-in-one baseline.
DrivingGaussian~\citep{wang2024drivinggaussian} further addresses large-scale
urban environments by compositing incremental static and dynamic Gaussian fields
with a composite dynamic Gaussian graph for multi-camera consistency.

\paragraph{Generative world models for autonomous driving.}
DriveDreamer~\citep{wang2023drivedreamer} is the first work to cast
autonomous driving scene generation explicitly as a world model problem,
training a diffusion model conditioned on structured traffic annotations and
ego-motion to produce photorealistic, temporally coherent multi-camera video.
DriveDreamer-2~\citep{zhao2024drivedreamer2} scales this approach with
LLM-generated traffic layouts, enabling greater diversity and user-controllable
scenario editing without re-training.
Drive-WM~\citep{wang2023drive_wm} proposes a multiview visual forecasting
architecture that couples scene prediction with planning, demonstrating
that a world model trained on real driving data can directly improve downstream
planning metrics.
OccWorld~\citep{zheng2024occworld} takes a 3-D occupancy perspective,
learning a world model that predicts future occupancy grids jointly with
ego-motion, unifying scene forecasting and motion planning in a single
autoregressive framework.
WoVoGen~\citep{lu2023wovogen} introduces a world-volume-aware diffusion model
that maintains spatial coherence across all cameras by reasoning over a
shared 3-D world volume rather than independent per-camera feature maps.

\paragraph{Diffusion-based street generation.}
StreetCrafter~\citep{yan2024streetcrafter} uses a point-cloud-conditioned video
diffusion model with controllable camera motion, enabling photorealistic
street-view synthesis for data augmentation.
Sat2Vid~\citep{wu2022sat2vid} lifts a satellite image to a street-view
panoramic video, addressing the altitude gap between aerial mapping and
ground-level perception.
Wild-GS~\citep{xu2024wildgs} solves the unconstrained internet-photo case by
disentangling transient and persistent Gaussian appearance in 3DGS.
InSpatio-WorldFM~\citep{chen2025inspatioworld} is an open-source real-time
generative frame model targeting spatial world understanding at scale.

\subsection{Generative World Models and Foundation Models}

World foundation models represent the apex of the NVS-to-world-modeling
trajectory: systems that jointly handle viewpoint, time, and action, trained
at internet scale.
We identify three design axes that differentiate current approaches.

\paragraph{Sequence vs.\ spatial representation.}
GAIA-1~\citep{hu2023gaia} treats driving scenes as sequences of video, text,
and action tokens, training an autoregressive transformer to predict the next
token.
This approach achieves strong temporal coherence but lacks an explicit 3-D
representation, limiting viewpoint flexibility.
Genie~\citep{bruce2024genie} learns a latent action space from unsupervised
internet videos, enabling interactive game-world generation without action labels---an
important step toward self-supervised world model pretraining.
WorldDreamer~\citep{worlddreamer2024} adopts a masked-token prediction
objective over a joint video-text tokenisation, demonstrating that masked
world model pre-training transfers well to diverse downstream generation tasks.

\paragraph{3-D-grounded generation.}
Genie~2~\citep{parkerholder2024genie2} introduces persistent 3-D structure
into the Genie framework, generating scenes that maintain geometry across
viewpoints and support first-person interactive navigation---the closest
current system to an NVS-grounded world model for open-ended environments.
WonderWorld~\citep{yu2024wonderworld} takes a complementary approach:
instead of learning from large-scale video, it uses instant 3DGS initialisation
from a single image combined with diffusion inpainting to generate interactive
3-D scenes on the fly.
WonderJourney~\citep{yu2023wonderjourney} chains text-guided scene generation
steps to produce perpetual explorable 3-D environments.
AVID~\citep{rigter2024avid} adapts pre-trained video diffusion models to
serve as general-purpose world models for offline reinforcement learning,
demonstrating that video generation and world modeling are more deeply
connected than previously thought.
Structured World Models~\citep{mendonca2023structured} learn
structured representations from human videos in a fully unsupervised manner,
enabling policy learning from video without any action labels.

\paragraph{Physical-world foundation models.}
Cosmos~\citep{agarwal2025cosmos} is the most ambitious current effort:
a family of world foundation models pretrained on internet-scale video with
explicit physical world conditioning, designed for downstream fine-tuning
to robotics and autonomous driving.
Pandora~\citep{xiang2024pandora} interleaves language and video tokens,
enabling language-driven world simulation with both visual and linguistic output.
The \emph{Is Sora a World Simulator?} survey~\citep{zhu2024sora_survey}
provides a comprehensive analysis of the relationship between video generation
models and world simulation capability, identifying key gaps in physical
consistency, causal understanding, and controllability.

\paragraph{The NVS-to-WM transition.}
A key architectural trend is the convergence of NVS and WM techniques:
whereas early world models like GAIA-1 operated in 2-D video space,
recent systems increasingly incorporate explicit 3-D representations---
whether Gaussian splats, triplane NeRFs, or occupancy grids---as the
``world state'' that the generative model predicts and renders.
This convergence has both practical and theoretical significance:
(i) 3-D-aware world models naturally support multi-camera and novel-viewpoint
prediction, which is essential for robotic and driving applications;
(ii) the NVS community's work on generalizable reconstruction provides the
initialization and geometry backbone that world models need;
and (iii) the world model community's work on temporal prediction and
action conditioning provides the dynamics prior that pure NVS systems lack.
The dual axis of \emph{spatial generalization} (D1--D2) and
\emph{temporal/action conditioning} (D3--D5, see Table~\ref{tab:nvs_wm})
captures this convergence precisely.

Table~\ref{tab:worldmodels} compares world modeling methods.

\begin{table}[t]
  \centering
  \caption{World modeling and neural scene simulation methods.}
  \label{tab:worldmodels}
  \scriptsize
  \setlength{\tabcolsep}{3pt}
  \begin{tabular}{@{}llccccc@{}}
    \toprule
    \textbf{Method} & \textbf{Year}
      & \textbf{3D Aware} & \textbf{Video} & \textbf{Action}
      & \textbf{Domain} & \textbf{Key Feature} \\
    \midrule
    GAIA-1~\citep{hu2023gaia}                  & 2023 & \pt  & \yes & \yes & Driving    & Seq. token WM \\
    UniSim~\citep{yang2023unisim}              & 2023 & \yes & \yes & \yes & Driving    & Sensor sim. \\
    MARS~\citep{wu2023mars}                    & 2023 & \yes & \yes & \pt  & Driving    & Compositional \\
    DriveDreamer~\citep{wang2023drivedreamer}  & 2024 & \pt  & \yes & \yes & Driving    & Diffusion WM \\
    OccWorld~\citep{zheng2024occworld}         & 2024 & \yes & \yes & \yes & Driving    & 3D occupancy WM \\
    WoVoGen~\citep{lu2023wovogen}              & 2024 & \yes & \yes & \yes & Driving    & World volume \\
    Drive-WM~\citep{wang2023drive_wm}          & 2024 & \pt  & \yes & \yes & Driving    & MV forecasting \\
    Genie~\citep{bruce2024genie}               & 2024 & \no  & \yes & \yes & Games      & Latent action \\
    Genie~2~\citep{parkerholder2024genie2}     & 2024 & \yes & \yes & \yes & Open-world & Interactive 3D \\
    WonderWorld~\citep{yu2024wonderworld}      & 2024 & \yes & \no  & \no  & Open-world & Instant 3DGS \\
    ViewCrafter~\citep{yu2024viewcrafter}      & 2024 & \yes & \yes & \no  & Open-world & Cam-cond video \\
    WorldDreamer~\citep{worlddreamer2024}      & 2024 & \no  & \yes & \pt  & General    & Masked token WM \\
    AVID~\citep{rigter2024avid}                & 2024 & \no  & \yes & \yes & General    & VDM$\to$WM \\
    Cosmos~\citep{agarwal2025cosmos}           & 2025 & \yes & \yes & \yes & Physical   & WFM training \\
    Street Gaussians~\citep{yan2024streetgaussian}& 2024& \yes & \yes & \pt & Driving   & 3DGS driving \\
    OmniRe~\citep{chen2024omnire}              & 2024 & \yes & \yes & \yes & Driving    & All-in-one \\
    DrivingGaussian~\citep{wang2024drivinggaussian}& 2024& \yes& \yes& \pt & Driving   & Large-scale GS \\
    \bottomrule
  \end{tabular}
\end{table}

\subsection{Embodied AI and Robotics}

NVS is more than a rendering tool in robotics: it provides the
\emph{spatial prior} that allows agents to reason about unseen parts of the
environment before taking action.
Three use patterns have emerged.

\textbf{Pose estimation and object localisation.}
iNeRF~\citep{yen2021inerf} inverts a trained NeRF to estimate 6-DoF camera
pose by gradient descent on the rendered photometric error---the first
demonstration that NeRF's differentiable renderer can be repurposed for
inference as well as reconstruction.
Dex-NeRF~\citep{ichnowski2021dex} applies NeRF to transparent-object grasping,
where conventional depth sensors fail, by leveraging NeRF's ability to
model view-dependent appearance to infer geometry from specular reflections.

\textbf{Visual planning and manipulation.}
NeRF-Supervisor~\citep{ze2023nerf_supervisor} uses a pre-trained NeRF to
synthesise robot-wrist camera views at hypothetical future poses, enabling
visual planning without physical rollouts.
GaussianGrasping~\citep{shorinwa2024gaussiangrasping} replaces NeRF with 3DGS
for real-time manipulation planning, exploiting the explicit Gaussian representation
for fast collision checking and grasp pose optimisation.

\textbf{Simulation-to-real transfer.}
RoboGSim~\citep{li2024robogssim} reconstructs robot workspaces in 3DGS and
renders synthetic training demonstrations from novel viewpoints, directly
closing the sim-to-real gap for visuomotor policy learning.

\subsection{Human Novel View Synthesis}

Synthesising novel views of humans is a critical subproblem with applications
in AR/VR, telepresence, and embodied AI.

\paragraph{Neural implicit human models.}
Neural Body~\citep{peng2021neuralbody} anchors structured latent codes on a
parametric human body (SMPL), enabling novel view and pose synthesis.
Rotationally-Temporally Consistent~\citep{liu2021rtcnvs} enforces both
rotation and temporal consistency for human performance video synthesis.
NHR~\citep{wu2020nhr} adopts multi-view point-cloud-based rendering for
dynamic human synthesis.

\paragraph{Generalizable and diffusion-based human NVS.}
CFSynthesis~\citep{li2023cfsynth} enables controllable and free-view 3-D human
video synthesis via a disentangled representation.
iVS-Net~\citep{chen2024ivsnet} learns human view synthesis from internet
videos without any 3-D supervision.
Blur2Sharp~\citep{dong2024blur2sharp} addresses human novel pose and view
synthesis from blurry input with a generative prior.
Diffuman4D~\citep{zheng2024diffuman4d} achieves 4-D consistent human view
synthesis from sparse-view videos via video diffusion.
EVA-Gaussian~\citep{hu2024evacg} achieves real-time 3DGS-based human NVS
under diverse multi-view settings.

\subsection{The Boundary: NVS vs.\ World Modeling}

Table~\ref{tab:nvs_wm} delineates the capability boundary.

\begin{table}[t]
  \centering
  \caption{NVS vs.\ World Modeling capability boundary.}
  \label{tab:nvs_wm}
  \small
  \setlength{\tabcolsep}{5pt}
  \begin{tabular}{@{}lcc@{}}
    \toprule
    \textbf{Capability} & \textbf{NVS (static)} & \textbf{World Model} \\
    \midrule
    Viewpoint-consistent rendering~(D1)   & \yes & \yes \\
    Sparse/single-image input~(D2)        & \pt  & \yes \\
    Temporal dynamics~(D3)                & \no  & \yes \\
    Action/agent conditioning~(D4)        & \no  & \yes \\
    Physical plausibility~(D5)            & \no  & \yes \\
    Cross-scene generalisation            & \pt  & \yes \\
    Causal scene understanding            & \no  & \yes \\
    \bottomrule
  \end{tabular}
\end{table}

% ─────────────────────────────────────────────────────────────────────────────
